% !TEX program = xelatex
% Lernplatt - by Can Siebert
% Kurs 71 (Dig 42) - Tag 8 - Loesungsheft

\documentclass[11pt,a4paper,oneside]{article}

\usepackage{polyglossia}
\setdefaultlanguage{german}
\usepackage[a4paper,margin=2.5cm]{geometry}

\usepackage{fontspec}
\defaultfontfeatures{Ligatures=TeX}
\IfFontExistsTF{Charter}{\setmainfont{Charter}}{%
  \IfFontExistsTF{Bitstream Charter}{\setmainfont{Bitstream Charter}}{\setmainfont{TeX Gyre Schola}}%
}
\IfFontExistsTF{Source Sans 3}{\setsansfont{Source Sans 3}}{%
  \IfFontExistsTF{Source Sans Pro}{\setsansfont{Source Sans Pro}}{\setsansfont{TeX Gyre Heros}}%
}
\newcommand{\caveatfontfallback}{\itshape}
\IfFontExistsTF{Caveat}{\newfontfamily\caveatfont{Caveat}[Scale=1.15]}{\newcommand\caveatfont{\caveatfontfallback}}
\setmonofont{Latin Modern Mono}

\usepackage{microtype}
\linespread{1.15}

\usepackage{xcolor}
\definecolor{lpInk}{RGB}{20,28,38}
\definecolor{lpAccent}{RGB}{157,74,26}
\definecolor{lpAccentLight}{RGB}{246,228,212}
\definecolor{lpAmber}{RGB}{222,138,30}
\definecolor{lpAmberLight}{RGB}{255,243,217}
\definecolor{lpAmberBorder}{RGB}{196,118,18}
\definecolor{lpGray}{RGB}{96,104,114}
\definecolor{lpGrayLight}{RGB}{238,240,243}
\definecolor{lpGreen}{RGB}{34,120,72}
\definecolor{lpGreenLight}{RGB}{222,238,228}
\definecolor{lpGreenBorder}{RGB}{24,96,58}
\definecolor{lpL1}{RGB}{34,120,72}
\definecolor{lpL2}{RGB}{22,90,140}
\definecolor{lpL3}{RGB}{170,42,52}

\usepackage[most]{tcolorbox}
\tcbuselibrary{breakable,skins}

\newtcolorbox{infobox}[1][Hinweis]{%
  enhanced, breakable, colback=lpAccentLight, colframe=lpAccent,
  boxrule=0.6pt, arc=2pt, left=10pt,right=10pt,top=8pt,bottom=8pt,
  fonttitle=\sffamily\bfseries\color{white}, coltitle=white,
  title={#1}, attach boxed title to top left={xshift=8pt,yshift=-8pt},
  boxed title style={size=small, colback=lpAccent, colframe=lpAccent, boxrule=0.4pt, arc=1pt},
  top=14pt
}

\newtcolorbox{loesungbox}[1][Musterl\"osung]{%
  enhanced, breakable, colback=lpGreenLight, colframe=lpGreenBorder,
  boxrule=0.6pt, arc=2pt, left=10pt,right=10pt,top=8pt,bottom=8pt,
  fonttitle=\sffamily\bfseries\color{white}, coltitle=white,
  title={#1}, attach boxed title to top left={xshift=8pt,yshift=-8pt},
  boxed title style={size=small, colback=lpGreenBorder, colframe=lpGreenBorder, boxrule=0.4pt, arc=1pt},
  top=14pt
}

\usepackage{enumitem}
\setlist{itemsep=2pt, topsep=4pt, parsep=0pt}
\setlist[itemize,1]{label={\textbullet},leftmargin=1.6em}
\setlist[itemize,2]{label={\textendash},leftmargin=1.4em}
\setlist[enumerate,1]{label=\arabic*., leftmargin=1.8em}

\usepackage{tabularx}
\usepackage{array}
\usepackage{booktabs}
\usepackage{longtable}

\usepackage{fancyhdr}
\usepackage{lastpage}
\pagestyle{fancy}
\fancyhf{}
\renewcommand{\headrulewidth}{0.3pt}
\renewcommand{\footrulewidth}{0pt}
\fancyhead[L]{\footnotesize\sffamily\color{lpGray}L\"osungsheft \textperiodcentered{} Kurs 71 \textperiodcentered{} Tag 8}
\fancyhead[R]{\footnotesize\sffamily\color{lpGray}Predictive \textperiodcentered{} DWH \textperiodcentered{} BI}
\fancyfoot[C]{\footnotesize\sffamily\color{lpGray}\thepage{} / \pageref{LastPage}}

\fancypagestyle{plain}{%
  \fancyhf{}%
  \renewcommand{\headrulewidth}{0pt}%
  \fancyfoot[C]{\footnotesize\sffamily\color{lpGray}\thepage{} / \pageref{LastPage}}%
}

\usepackage[explicit]{titlesec}
\titleformat{\section}[block]
  {\sffamily\Large\bfseries\color{lpAccent}}
  {\thesection.}{0.6em}{#1}
  [{\vspace{2pt}\color{lpAccent}\titlerule[0.6pt]}]
\titleformat{\subsection}[block]
  {\sffamily\large\bfseries\color{lpInk}}
  {\thesubsection}{0.6em}{#1}
\titlespacing*{\section}{0pt}{14pt}{8pt}
\titlespacing*{\subsection}{0pt}{10pt}{4pt}

\usepackage[hidelinks,unicode]{hyperref}

\newcommand{\akzent}[1]{{\caveatfont\color{lpAmber}#1}}
\newcommand{\fachbegriff}[1]{\textbf{#1}}
\newcommand{\quelle}[1]{{\footnotesize\sffamily\color{lpGray}(#1)}}

\newcommand{\niveaubadge}[2]{%
  \tcbox[on line, colback=#1, colframe=#1, coltext=white,
         boxrule=0pt, arc=2pt, left=5pt,right=5pt,top=1pt,bottom=1pt,
         nobeforeafter]{\sffamily\bfseries\footnotesize #2}%
}
\newcommand{\nivLi}{\niveaubadge{lpL1}{L1\,\textbullet\,Wissen}}
\newcommand{\nivLii}{\niveaubadge{lpL2}{L2\,\textbullet\,Anwendung}}
\newcommand{\nivLiii}{\niveaubadge{lpL3}{L3\,\textbullet\,Management}}

\newcommand{\zeit}[1]{%
  \tcbox[on line, colback=lpGrayLight, colframe=lpGrayLight, coltext=lpInk,
         boxrule=0pt, arc=2pt, left=5pt,right=5pt,top=1pt,bottom=1pt,
         nobeforeafter]{\sffamily\footnotesize\textbf{$\sim$\,#1}}%
}

\newcommand{\aufgabenkopf}[4]{%
  \par\vspace{6pt}
  \noindent{\sffamily\Large\bfseries\color{lpAccent}Aufgabe #1 \textendash{} #2}\par
  \vspace{2pt}
  \noindent{\color{lpAccent}\rule{\linewidth}{0.6pt}}\par
  \vspace{4pt}
  \noindent #3 \hfill \zeit{#4}\par
  \vspace{6pt}
}

\newcommand{\ytquery}[1]{%
  \par\vspace{4pt}
  \noindent{\footnotesize\sffamily\color{lpGray}\textbf{YouTube-Suchquery:} \texttt{#1}}\par
}

\begin{document}

\begin{titlepage}
  \pagestyle{empty}
  \vspace*{1.5cm}
  \noindent{\sffamily\footnotesize\color{lpGray}LERNPLATT \textperiodcentered{} BY CAN SIEBERT}\\[2pt]
  \noindent{\color{lpAccent}\rule{\linewidth}{1.2pt}}\\[4pt]
  \noindent{\sffamily\footnotesize\color{lpGray}L\"OSUNGSHEFT \textperiodcentered{} MODUL 3 \textperiodcentered{} TAG 8 VON 16 \textperiodcentered{} KURS 71 (Dig 42) \textperiodcentered{} 9.\,Juni 2026}

  \vspace{3cm}
  {\sffamily\fontsize{30}{34}\selectfont\bfseries\color{lpInk}L\"osungsheft\\Tag 8\par}

  \vspace{0.7cm}
  {\sffamily\large\color{lpGray}Musterl\"osungen \textendash{} Enhancement,\\Predictive, Real-time \& BI.\par}

  \vspace{1.5cm}
  {\caveatfont\LARGE\color{lpGreen}Musterl\"osungen \textperiodcentered{} nicht die einzige Antwort}\par

  \vfill

  \noindent\begin{minipage}{0.55\linewidth}
    {\sffamily\footnotesize\color{lpGray}Hinweis:}\\
    {\sffamily\small\color{lpInk}Musterl\"osungen sind Diskussionsbasis,\\nicht die einzige korrekte L\"osung.}\\[10pt]
    {\sffamily\footnotesize\color{lpGray}Begleitend:}\\
    {\sffamily\small\color{lpInk}Aufgabenheft \& Lehrskript Tag 8}
  \end{minipage}\hfill
  \begin{minipage}{0.4\linewidth}
    \raggedleft
    {\sffamily\footnotesize\color{lpGray}Dozent:}\\
    {\sffamily\small\color{lpInk}Can Siebert}\\[10pt]
    {\sffamily\footnotesize\color{lpGray}Niveau-Mix:}\\
    {\sffamily\small\color{lpInk}3\,L1 \textperiodcentered{} 5\,L2 \textperiodcentered{} 3\,L3}
  \end{minipage}

  \vspace{0.5cm}
  \noindent{\color{lpAccent}\rule{\linewidth}{0.6pt}}
\end{titlepage}

\section*{Hinweise zu den Musterl\"osungen}
\thispagestyle{plain}

\noindent Heute steht der Senior-Reflex im Vordergrund: \emph{Predictive ohne Handlungsplan ist wertlos}. Alert \textrightarrow{} Owner \textrightarrow{} Playbook. Und: \emph{Baseline vor ML} \textendash{} einfache Regeln zuerst.

\begin{infobox}[Nutzung im Coaching]
Pr\"ufen Sie Ihre L\"osung gegen drei Fragen: (1) Habe ich die Baseline-Regel ernsthaft als Alternative gepr\"uft? (2) Hat jeder Alert einen Owner und ein Playbook? (3) Habe ich Drift-Monitoring eingebaut?
\end{infobox}

\clearpage

\section*{Niveau L1 \textendash{} Wissen \& Reproduktion}
\addcontentsline{toc}{section}{L1 -- Wissen}

% LERNPLATT_HILFSVIDEOS
\section*{Hilfsvideos zu diesem Tag}
\addcontentsline{toc}{section}{Hilfsvideos zu diesem Tag}
\thispagestyle{plain}

\noindent Die folgenden YouTube-Erkl\"arvideos vermitteln die Inhalte, die Sie zur Bearbeitung der Aufgaben ben\"otigen. Klicken Sie auf den Titel, um das Video direkt zu \"offnen; die vollst\"andige URL steht in Klammern.

\begin{description}[style=nextline,leftmargin=18pt,labelindent=0pt,itemsep=8pt]
  \item[1.~ETL (Extract Transform Load) für Data Warehousing einfach erklärt] \href{https://youtu.be/VVJRBhUkX5o}{\textcolor{lpAccent}{https://youtu.be/VVJRBhUkX5o}}\\[2pt]
  {\footnotesize\color{lpGray}(URL: \nolinkurl{https://youtu.be/VVJRBhUkX5o})}
  \item[2.~Datenbank vs. Data Lake vs. Data Warehouse — die Unterschiede] \href{https://youtu.be/BQuyIwi_RCQ}{\textcolor{lpAccent}{https://youtu.be/BQuyIwi_RCQ}}\\[2pt]
  {\footnotesize\color{lpGray}(URL: \nolinkurl{https://youtu.be/BQuyIwi_RCQ})}
  \item[3.~Data Warehouse einfach erklärt — Grundlagen] \href{https://youtu.be/EHmeK1nMBvA}{\textcolor{lpAccent}{https://youtu.be/EHmeK1nMBvA}}\\[2pt]
  {\footnotesize\color{lpGray}(URL: \nolinkurl{https://youtu.be/EHmeK1nMBvA})}
  \item[4.~Business Intelligence KPI Dashboard in Power BI] \href{https://youtu.be/0BwQvLGUbZg}{\textcolor{lpAccent}{https://youtu.be/0BwQvLGUbZg}}\\[2pt]
  {\footnotesize\color{lpGray}(URL: \nolinkurl{https://youtu.be/0BwQvLGUbZg})}
\end{description}
\vspace{0.6em}
\noindent{\footnotesize\color{lpGray}\textit{Tipp:} \"Offnen Sie das Video, lassen Sie es im Hintergrund laufen und arbeiten Sie parallel an der Aufgabe.}

\clearpage

\aufgabenkopf{1}{ML-Typen im Prozesskontext}{\nivLi}{8\,min}

\ytquery{Machine Learning Process Mining Classification Regression Clustering Anomaly}

\begin{loesungbox}
\begin{tabularx}{\linewidth}{@{}p{3.0cm}|X|X|X@{}}
\toprule
\textbf{Typ} & \textbf{Use Case} & \textbf{Datenminimum} & \textbf{Grenzen} \\
\midrule
Klassifikation & ``Eskalation ja/nein?''         & gelabelte F\"alle (positive + negative) & Datenqualit\"at; Label-Bias. \\
\midrule
Regression     & Restlaufzeit prognostizieren    & numerische Zielgr\"o{\ss}e + Historie       & Extreme; Trendbrueche. \\
\midrule
Clustering     & Falltypen erkennen              & repr\"asentative Stichprobe         & Cluster-Interpretation; Drift. \\
\midrule
Anomalie       & ungew\"ohnliche Pfade           & ausreichend ``normale'' Cases     & False Positives; Erkl\"arbarkeit. \\
\bottomrule
\end{tabularx}

\smallskip
\emph{Allgemeine Grenzen:} Konzeptdrift (Welt \"andert sich), Fehlanreize (Modell triggert Verhalten), Erkl\"arbarkeit (Black-Box vs.\ Audit), und nicht zuletzt: Datenqualit\"at als Grundproblem.
\end{loesungbox}

\clearpage

\aufgabenkopf{2}{ETL-Pipeline}{\nivLi}{8\,min}

\ytquery{ETL Pipeline Process Data Staging Mapping Join Data Quality DWH}

\begin{loesungbox}
\begin{tabularx}{\linewidth}{@{}p{2.4cm}XX@{}}
\toprule
\textbf{Schritt} & \textbf{Was passiert?} & \textbf{Stolperstein} \\
\midrule
Quelle         & CRM/ERP/Tickets exportieren. & Schnittstellen-\"Anderung ohne Notice. \\
Staging        & Rohdaten unver\"andert ablegen. & Disk-Volumen, Aufbewahrungsfristen. \\
Mapping        & Aktivit\"aten / Status normalisieren. & Inkonsistente Synonyme \"ubersehen. \\
Join           & \"uber Case-ID / Korrelations\-schl\"ussel. & ID-Kollisionen, missende Joins. \\
Quality Checks & Pflichtfelder, Zeit, Dubletten. & Checks zu lasch (laufen alle durch). \\
DWH            & Faktentabelle Events + Dimensionen. & Modell zu eng (sp\"atere KPIs unm\"oglich). \\
\bottomrule
\end{tabularx}

\smallskip
\textbf{Meiste Arbeit:} Mapping + Quality Checks (oft 60\,\% des Aufwands). \textbf{Gr\"o{\ss}ter Wert:} Stabilit\"at + Konsistenz im DWH \textendash{} alles dr\"uber baut darauf auf.
\end{loesungbox}

\clearpage

\aufgabenkopf{3}{BI-Dashboard}{\nivLi}{10\,min}

\ytquery{BI Dashboard Best Practice Executive KPI Drilldown Alert Tableau Power BI}

\begin{loesungbox}
\textbf{Drei Prinzipien:}
\begin{enumerate}
  \item \emph{Wenige Executive-KPIs:} max.\ 4\,--\,6, sonst \"uberfordert. Anti-Muster (``zuviel''): Kennzahlen-Tapete, niemand handelt.
  \item \emph{Drilldown:} ein Klick auf KPI zeigt Ursache. Anti-Muster (``zuwenig''): KPI als Standalone-Zahl, Fragen offen.
  \item \emph{Alarme:} Schwellenwert + Owner + n\"achster Schritt. Anti-Muster (``zuwenig''): Anzeige ohne Handlungsaufforderung.
\end{enumerate}

\smallskip
\textbf{Mini-Layout ``Incident Operations'':}
\begin{enumerate}
  \item \emph{Exec KPI 1:} SLA-Einhaltungsquote (Trend 30\,T).
  \item \emph{Exec KPI 2:} Mean Time to Resolution.
  \item \emph{Prozess-KPI 3:} Top-10 Risikotickets (Drilldown).
  \item \emph{Prozess-KPI 4:} Reassignment-Quote pro Service.
  \item \emph{Risk/Compliance:} Anzahl audit-relevanter Incidents.
  \item \emph{Alert:} Live-Liste der Tickets mit Risk-Score $>$\,0,8 (Owner sichtbar).
\end{enumerate}
\end{loesungbox}

\clearpage

\section*{Niveau L2 \textendash{} Anwendung \& Transfer}
\addcontentsline{toc}{section}{L2 -- Anwendung}

\aufgabenkopf{4}{Predictive Use Case Canvas (zwei Use Cases)}{\nivLii}{25\,min}

\ytquery{Predictive Process Analytics Use Case Canvas SLA Risk Rework Baseline}

\begin{loesungbox}
\textbf{Use Case A \textendash{} SLA-Risiko fr\"uh erkennen:}
\begin{itemize}
  \item Zielentscheidung: Ressourcen umschichten bevor SLA bricht.
  \item Prognose: Wahrscheinlichkeit SLA-Breach in n\"achsten 4\,h.
  \item Handlung: Duty Manager bekommt Top-10-Liste, Triage 15\,Min, Ressourcen-Shift.
  \item Datenminimum: Priorit\"at, Ticket-Alter, Queue, Reassignments, Service.
  \item Risiko/Compliance: keine Bias gegen Service-A/B; Audit-Log aller Alerts.
  \item Erfolg: SLA-Breach-Quote sinkt um 20\,\% in 8\,Wochen.
\end{itemize}

\smallskip
\textbf{Use Case B \textendash{} Rework / Fehler fr\"uh erkennen:}
\begin{itemize}
  \item Zielentscheidung: Coaching f\"ur Team / DoR sch\"arfen.
  \item Prognose: Wahrscheinlichkeit Rework $>$\,1 Schleife.
  \item Handlung: Team-Lead bekommt Liste; Coaching-Session.
  \item Datenminimum: Auftragstyp, Bearbeiter, Eingangs-Kanal, Datenvollst\"andigkeit, fr\"uhere Rework-Quote.
  \item Risiko/Compliance: keine Sanktion einzelner MA, nur Team-Coaching; Audit-Log.
  \item Erfolg: Rework-Quote sinkt um 15\,\%.
\end{itemize}

\smallskip
\textbf{Baseline-Heuristiken:}
\begin{itemize}
  \item A: ``Ticket-Alter $>$\,70\,\% SLA \textrightarrow{} Risiko hoch''.
  \item B: ``$\geq$\,2 fehlende Pflichtfelder bei Eingang \textrightarrow{} Rework-Risiko hoch''.
\end{itemize}

\smallskip
\textbf{Fehlende Daten pragmatisch erfassen:}\\
\emph{Bearbeiter-ID bei Telefon-Tickets} fehlt h\"aufig. \emph{Pragmatischer Weg:} Pflichtfeld im Ticket-Tool, Default ``unknown'' nur f\"ur 3 Wochen, dann verpflichtend.
\end{loesungbox}

\clearpage

\aufgabenkopf{5}{ETL \& Dashboard in 30 Minuten}{\nivLii}{25\,min}

\ytquery{ETL Pipeline 6 Steps Data Quality Checks BI Dashboard Anti Gaming}

\begin{loesungbox}
\textbf{ETL-Pipeline (6 Schritte):}\\
1.\ CRM/Ticketing/ERP exportieren \textrightarrow{} 2.\ Staging-Schicht \textrightarrow{} 3.\ Aktivit\"atsname-Mapping (40\,Synonyme \textrightarrow{} 12\,Standards) \textrightarrow{} 4.\ Join \"uber Customer-Key + Order-ID \textrightarrow{} 5.\ Quality Checks (s.\,u.) \textrightarrow{} 6.\ DWH ``Faktentabelle Events'' + Dim (Service, Kunde, Zeit, Ressource).

\smallskip
\textbf{F\"unf Datenqualit\"ats-Checks:}
\begin{enumerate}
  \item Fehlende Case-ID (\texttt{NULL}).
  \item Zeitstempel r\"uckw\"arts ($\Delta t < 0$).
  \item Unbekannte Aktivit\"at (nicht im Mapping).
  \item Dublette (gleiche Case-ID + Aktivit\"at + Timestamp $\pm$ 1\,s).
  \item Leeres Pflichtfeld (Owner-ID).
\end{enumerate}

\smallskip
\textbf{BI-Dashboard mit 6 Kacheln:}
\begin{enumerate}
  \item Exec: SLA-Einhaltung 30\,T-Trend.
  \item Exec: MTTR (mean time to resolution).
  \item Prozess: Reassignment-Rate pro Service.
  \item Prozess: Top-10 Risikotickets.
  \item Risk: Anzahl audit-relevanter Incidents.
  \item Alert: Live-Risikoliste mit Owner + ``Next Best Action''.
\end{enumerate}

\smallskip
\textbf{Anti-Gaming-Entscheidung:}\\
\emph{``Anzahl gel\"oster Tickets pro MA pro Tag''} f\"uhren wir nicht ein. Begr\"undung: belohnt Splitting; KPI besser \"uber Service-Level (Team-Performance), nicht Individuum.
\end{loesungbox}

\clearpage

\aufgabenkopf{6}{Baseline-Regel vor ML}{\nivLii}{18\,min}

\ytquery{Baseline Heuristic vs Machine Learning When Simpler is Better Process}

\begin{loesungbox}
\textbf{1. Baseline-Regel ``SLA-Breach in 4\,h'':}\\
\emph{Risk = hoch wenn Ticket-Alter $>$\,70\,\% der SLA-Zeit \emph{oder} (Reassignments $\geq$ 2 \emph{und} Ticket-Alter $>$\,50\,\% SLA).}

\smallskip
\textbf{2. Messbarkeit:}
\begin{itemize}
  \item KPI A: \emph{Precision} (von den Alerts \textendash{} wie viele waren echte Breaches?).
  \item KPI B: \emph{Recall} (von den echten Breaches \textendash{} wie viele wurden alarmiert?).
\end{itemize}
Vergleich vor Baseline (``manuell'') vs.\ mit Baseline vs.\ sp\"ater ML.

\smallskip
\textbf{3. Wann lohnt ML?}
\begin{itemize}
  \item Baseline-Precision dauerhaft $<$\,60\,\% (zu viele False Positives).
  \item Baseline-Recall $<$\,70\,\% (zu viele verpasste).
  \item Datenvolumen + Features ausreichend f\"ur Training, und der Aufwand f\"ur MLOps lohnt sich (Monitoring, Drift).
\end{itemize}

\smallskip
\textbf{4. Anti-Muster: Baseline besser als ML?}\\
Bei \emph{regulierten Prozessen} (Pharma, Finance) ist eine Regel auditierbar und stabil \textendash{} ein ML-Modell muss erkl\"arbar gemacht werden und ist drift-anf\"allig. F\"ur sehr seltene Ereignisse (kleines Label-Volumen) liefert ML wenig Mehrwert.
\end{loesungbox}

\clearpage

\aufgabenkopf{7}{Alert-Playbook}{\nivLii}{20\,min}

\ytquery{Alert Playbook Process Mining MLOps Drift Detection Alarm Fatigue}

\begin{loesungbox}
\textbf{1. Alert-Playbook ``SLA-Risiko'':}

\begin{tabularx}{\linewidth}{@{}p{2.4cm}|X|l|X@{}}
\toprule
\textbf{Stufe} & \textbf{Owner} & \textbf{Handlung} & \textbf{Eskalation/Doc} \\
\midrule
Risk 0,5\,--\,0,7 & Team-Lead & Pr\"ufen + ggf.\ Reassign & Im Ticket vermerken. \\
Risk 0,7\,--\,0,85 & Duty Manager & Fast Triage 15\,Min + Ressourcen-Shift & Wenn nicht gel\"ost: Eskalation an SDM. \\
Risk $>$\,0,85 & Service Delivery Manager & Sofortma{\ss}nahme + Kundeninfo & Post-Review verpflichtend. \\
\bottomrule
\end{tabularx}

\smallskip
\textbf{2. Alarmm\"udigkeit verhindern:}
\begin{itemize}
  \item \emph{Schwellen} dynamisch je Service.
  \item \emph{Eskalations\-pfad}: nicht jeder Alert geht an alle \textendash{} nur Stufe drei trifft den SDM.
  \item \emph{Feedback-Schleife}: jede Woche werden 5 Alerts manuell \"uberpr\"uft (``war Alert n\"utzlich?''), Schwellen werden danach justiert.
\end{itemize}

\smallskip
\textbf{3. Drift-Monitoring:}
\begin{itemize}
  \item \emph{Data Drift}: Feature-Verteilung verschiebt sich (z.\,B.\ ``Ticket-Alter''-Median +30\,\%) \textrightarrow{} Pr\"ufung Quellsystem.
  \item \emph{Concept Drift}: Precision/Recall sinken bei stabiler Feature-Verteilung \textrightarrow{} Welt hat sich ge\"andert; Retraining oder Regel-Update.
\end{itemize}

\smallskip
\textbf{4. Faustregel:} Ein Predictive-System ist \emph{produktiv}, sobald (a) die Quote ``Alerts mit dokumentierter Handlung'' $>$\,80\,\% liegt, (b) Drift-Monitoring l\"auft, (c) ein dedizierter Owner pro Modell vorhanden ist.
\end{loesungbox}

\clearpage

\aufgabenkopf{8}{False Positive vs.\ False Negative}{\nivLii}{18\,min}

\ytquery{False Positive vs False Negative Cost Sensitive ML Process Mining Threshold}

\begin{loesungbox}
\begin{tabularx}{\linewidth}{@{}p{2.4cm}XXl@{}}
\toprule
\textbf{Use Case} & \textbf{False Positive} & \textbf{False Negative} & \textbf{Teurer} \\
\midrule
SLA-Breach    & Alarm ohne Breach (Team gestresst, Kunden unsicher). & Verpasster Breach (Vertragsstrafe, Kunde unzufrieden). & FN $\gg$ FP \\
\midrule
Maverick Buying & Buchung verz\"ogert (Reklamation Einkauf). & Compliance-Versto{\ss} (Bu{\ss}geld). & FN $\gg$ FP \\
\midrule
Kunde springt ab & Unn\"otiger Anruf (Aufwand Customer Success). & Kunde springt wirklich ab (Umsatz verloren). & FN $>$ FP \\
\midrule
Versand-Verz\"ogerung & Vorwarnung ohne Realit\"at (Kunde irritiert). & Tats\"achliche Verz\"ogerung (Kunde sauer + Storno). & FN $>$ FP \\
\bottomrule
\end{tabularx}

\smallskip
\textbf{Konsequenz f\"ur den Schwellenwert:}\\
Wenn \emph{False Negative} teurer ist, w\"ahlt man \emph{niedrige Schwelle} (sensitiv) \textendash{} lieber ein paar Fehlalarme als verpasste F\"alle. Konkret f\"ur SLA: Risk-Schwelle 0,5\,--\,0,6 statt 0,8.

\emph{Caveat:} Zu niedrige Schwelle erzeugt Alarmm\"udigkeit \textendash{} dann kehrt sich das Problem um. Schwelle wird daher iterativ gegen das Alarmm\"udigkeits-Signal justiert.
\end{loesungbox}

\clearpage

\section*{Niveau L3 \textendash{} Management \& Entscheidungsarchitektur}
\addcontentsline{toc}{section}{L3 -- Management}

\aufgabenkopf{9}{Fallstudie \emph{FinServe IT Operations}}{\nivLiii}{45\,min}

\ytquery{SLA Predictive Incident Management Features ETL Dashboard Real Time}

\begin{loesungbox}
\textbf{1. Predictive Use Case:} ``Wahrscheinlichkeit SLA-Breach in n\"achsten 4\,h''.

\smallskip
\textbf{2. Features (10):} Priority, Ticket Age, Queue, Reassignments, Service, Category, Open Incidents per Service, Time since last update, Shift, Customer Tier.

\smallskip
\textbf{3. Label-Definition:} \emph{SLA-Breach} = Ticket nicht gel\"ost vor SLA-Zeit. \emph{Fehlende Labels} (Tickets noch offen): Proxy-Regel ``Ticket-Age $>$\,70\,\% SLA'' f\"ur Training.

\smallskip
\textbf{4. ETL-Pipeline:} Ticketing-Extract \textrightarrow{} Staging \textrightarrow{} Mapping (Service/Kategorie) \textrightarrow{} Join Monitoring-Signals \textrightarrow{} Quality Checks (Case-ID, Timestamps, Dubletten) \textrightarrow{} DWH ``Events'' + ``Ticket Facts''. \emph{Case-ID-Strategie:} Ticket-ID als Master; bei zusammengef\"uhrten Tickets parent-child-Relation im DWH abbilden.

\smallskip
\textbf{5. Dashboard:}
\begin{itemize}
  \item SLA-Risk Heatmap (Service $\times$ Queue).
  \item Top-10 Risikotickets.
  \item Trend Breach-Rate (4\,W).
  \item Avg.\ Resolution Time.
  \item Reassignments pro Service.
  \item Alert-Log mit Owner.
\end{itemize}

\smallskip
\textbf{Playbook:} Risk $>$\,0,8 \textrightarrow{} Duty Manager notified \textrightarrow{} Fast Triage 15\,Min \textrightarrow{} Ressourcen umschichten \textrightarrow{} Dokumentation im Ticket \textrightarrow{} Post-Review bei wiederkehrenden Mustern.

\smallskip
\textbf{6. MVP-Entscheidung:} \emph{Weglassen:} (a) Monitoring-Join (eigener Stream zu fragil), (b) Customer-Tier (Datenschutz-Pr\"ufung l\"auft). \emph{Kompensieren:} manuelle ``Service Impact''-Kennzeichnung im Ticket; Service-Tier statt Kunde-Tier.
\end{loesungbox}

\clearpage

\aufgabenkopf{10}{Modell-Governance}{\nivLiii}{25\,min}

\ytquery{MLOps Model Governance Drift Detection Explainability XAI Fairness Bias}

\begin{loesungbox}
\textbf{1. ML Product Owner (SLA-Modell):}
\begin{itemize}
  \item Aufgaben: Use Case definieren, Modell-Lifecycle (Train/Retrain/Decom), Monitoring, Stakeholder-Kommunikation.
  \item Mandat: darf Schwellen anpassen, darf Modell still legen.
  \item Eskalation: Owner \textrightarrow{} Head of Process Excellence \textrightarrow{} CDO bei Auswirkungen auf Compliance / mehrere Bereiche.
\end{itemize}

\smallskip
\textbf{2. Drift-Typen:}
\begin{itemize}
  \item \emph{Data Drift:} Feature-Verteilung verschiebt sich (z.\,B.\ neue Service-Kategorie). Messen: KL-Divergenz / PSI pro Feature.
  \item \emph{Concept Drift:} Beziehung Feature\,\textrightarrow{}\,Label \"andert sich. Messen: Precision/Recall sinkt bei stabiler Feature-Verteilung.
\end{itemize}

\smallskip
\textbf{3. Erkl\"arbarkeit:}
\begin{itemize}
  \item \emph{Feature Importance} (global): welche Features tragen am meisten?
  \item \emph{SHAP-Werte} (lokal): warum hat \emph{dieses} Ticket Risk 0,87? (Top-3 Treiber visualisieren.)
\end{itemize}
F\"ur Duty Manager: lokale Erkl\"arung pro Alert (``Risiko wegen: hohe Reassignments, Queue lang, Tageszeit'').

\smallskip
\textbf{4. Fairness / Bias:}\\
Modell k\"onnte Service-A systematisch hochstufen, weil dort historisch viele Breaches waren. Folge: Service-A wird ``\"uberbetreut'', Service-B unter. \emph{Gegenma{\ss}nahme:} Service-stratifizierte Performance-Reports; bei systematischem Bias Re-Training mit Sampling-Strategie.
\end{loesungbox}

\clearpage

\aufgabenkopf{11}{Transferprojekt \textendash{} Process Intelligence Platform}{\nivLiii}{45\,min}

\ytquery{Process Intelligence Platform Decision Architecture CDO Real Time Predictive}

\begin{loesungbox}
\textbf{1. Top-10 Entscheidungen:}\\
Priorisierung Use Cases \textperiodcentered{} SLA-Steuerung \textperiodcentered{} Kapazit\"aten \textperiodcentered{} Risiko (Compliance) \textperiodcentered{} Automatisierung \textperiodcentered{} Investitionen \textperiodcentered{} Pricing \textperiodcentered{} Workforce \textperiodcentered{} Audit-Befragung \textperiodcentered{} Vorstands-Reporting.

\smallskip
\textbf{2. Data Product Design:}
\begin{itemize}
  \item Event-Log-Standard: 5 Pflicht + 5 Optional-Felder.
  \item Case-ID-Strategie: Master-System + UUID.
  \item Datenmodell: Faktentabelle Events + Dim (Zeit, Kunde, Produkt, Standort, Ressource).
  \item DQ-SLAs: Vollst\"andigkeit $\geq$\,95\,\%, Aktualit\"at $\leq$\,4\,h.
  \item Ownership: Data Owner pro Quellsystem.
\end{itemize}

\smallskip
\textbf{3. Analytics Portfolio:} Value/Risk/Effort-Matrix, Baseline-First-Policy, MLOps mit Drift-Monitoring, Erkl\"arbarkeit per default.

\smallskip
\textbf{4. Real-time Operating Model:}
\begin{itemize}
  \item Alerting Governance: pro Use Case ein Playbook.
  \item Eskalations\-logik: Team-Lead \textrightarrow{} Duty Manager \textrightarrow{} SDM.
  \item Audit Trail: jeder Alert + Handlung dokumentiert, 7\,Jahre.
  \item Human-in-the-loop: bei kritischen Entscheidungen (z.\,B.\ Compliance) immer Mensch.
\end{itemize}

\smallskip
\textbf{5. BI \& KPI Governance:} Definition-Tag, monatlicher Review, Anti-Gaming-Kopplungen (siehe Tag 7).

\smallskip
\textbf{6. Roadmap (20\,W):}
\begin{itemize}
  \item Wo 1\,--\,8 (MVP): 1 Use Case (SLA), DWH-Grundger\"ust, 1 Dashboard, Baseline-Regel produktiv.
  \item Wo 9\,--\,16 (Scale): 2 weitere Use Cases, ML-Modell f\"ur SLA, Drift-Monitoring.
  \item Wo 17\,--\,20 (Sustain): Governance Board, Erkl\"arbarkeits-Layer, Audit-Vorbereitung.
\end{itemize}

\smallskip
\textbf{7. Entscheidungen unter Unsicherheit:}
\begin{enumerate}
  \item \emph{3 Use Cases zuerst:} SLA-Breach, Rework-Risiko, Maverick Buying. Verworfen: ``Vorhersage Kundenabsprung'' (Datenschutz-Aufwand zu hoch). Fr\"uhwarnsignal: wenn nach 8\,W kein Use Case in Produktion \textrightarrow{} Scope auf 1 Use Case reduzieren.
  \item \emph{Executive-KPIs:} SLA-Einhaltung + MTTR + Audit-Findings. \emph{Guardrails:} keine Pro-Mitarbeiter-Metriken; Definition-Tag; Anti-Gaming \"uber gekoppelte Kennzahlen.
\end{enumerate}
\end{loesungbox}

\clearpage

\section*{Abschluss}
\thispagestyle{plain}

\noindent Die Musterl\"osungen sind eine Diskussionsbasis. Heute pr\"ufen Sie: Hat jede Prognose einen Handlungsplan? Ist die Baseline ernsthaft als Alternative gepr\"uft? Habe ich Drift-Monitoring eingebaut?

\medskip
\begin{infobox}[N\"achster Schritt]
Bringen Sie ins Coaching: Welche \emph{eine} Baseline-Regel w\"urden Sie morgen aktivieren? Welcher Alert h\"atte heute einen Owner? Welche Erkl\"arung w\"urden Sie dem CFO geben, wenn das Modell drei Wochen lang ``schweigt''?
\end{infobox}

\vspace{1cm}
\noindent{\color{lpAccent}\rule{\linewidth}{0.6pt}}\\[4pt]
{\footnotesize\sffamily\color{lpGray}Lernplatt \textperiodcentered{} by Can Siebert \textperiodcentered{} Kurs 71 (Dig 42) \textperiodcentered{} Tag 8 \textperiodcentered{} L\"osungsheft \textperiodcentered{} \textcopyright{} 2026}

\end{document}
